\begin{abstract}
Multimodal Large Language Models (MLLMs) have demonstrated strong capabilities
in visual question answering and image-text reasoning,
yet they remain susceptible to hallucinations,
where generated outputs are inconsistent with visual inputs,
contradict established world knowledge,
or contain flawed reasoning.
This report presents a systematic framework
for detecting and evaluating hallucinations in MLLMs.
We propose a three-category taxonomy
covering faithfulness, factuality, and logical hallucinations,
and evaluate four representative models
(GPT-5.4-mini, Gemini~2.5~Flash, Qwen3.5-35B-A3B, Qwen3-VL-235B-A22B)
across three benchmarks:
POPE for object existence probing,
MathVista for mathematical visual reasoning,
and VQA-RAD for medical radiology VQA.
Two detection methods are employed:
a rule-based approach for binary probing tasks
and an MLLM Judge protocol based on the MMHal-Bench 0--6 Likert scale
for open-ended evaluation.
Key findings include:
(1)~model rankings shift substantially across task domains,
with no single model dominating all scenarios;
(2)~Chain-of-Thought prompting significantly reduces hallucinations
in weaker-baseline models
but introduces additional risks for stronger ones
through exposure effects and commitment errors;
and (3)~faithfulness hallucinations account for over 96\%
in medical VQA,
indicating that fine-grained visual perception
remains the primary bottleneck in specialized domains.
Comprehensive error analyses,
including human alignment verification (Cohen's $\kappa = 0.70$),
cross-judge consistency checks,
threshold sensitivity analysis,
and length bias assessment,
confirm the robustness of our evaluation framework.
\end{abstract}
